\documentclass[12pt]{article}
\usepackage[T1]{fontenc}
\usepackage[utf8]{inputenc}
\usepackage{lmodern}
\usepackage{textcomp}
\usepackage{enumitem}
\usepackage{geometry}
\geometry{margin=1in}
\begin{document}
\begin{center}
Answer Key - Geography - Undergraduate Level Assessment 19 \
\small (Answers are ordered exactly as in the question paper.)
\end{center}

\section*{Multiple Choice}
\begin{enumerate}[label=\arabic*.]
\item \textbf{Answer:} A
\\ \textit{Options:} A) China; B) United States; C) Germany; D) Japan
\\ \textit{Note:} China generated the most total energy from solar sources in 2019 due to its massive investments in solar technology, extensive manufacturing capabilities, and the implementation of supportive government policies that accelerated the deployment of solar energy systems.
\item \textbf{Answer:} B
\\ \textit{Options:} A) strongly supported by most studies; B) supported mainly by cross-section, not time-series studies; C) supported mainly by time-series, not cross-section studies; D) generally repudiated by empirical studies
\\ \textit{Note:} The correct answer highlights that the "inverted U hypothesis" regarding the relationship between inequality and development has primarily been evidenced through cross-sectional studies, which analyze data at a single point in time, rather than time-series studies that track changes over time, indicating a lack of longitudinal data to robustly support the hypothesis.
\item \textbf{Answer:} C
\\ \textit{Options:} A) 1\%; B) 2\%; C) 4\%; D) 8\%
\\ \textit{Note:} The correct answer is 4\% because, as of 2017, global public spending on education was estimated to represent approximately this percentage of total global GDP, reflecting the significant investment made by governments in educational systems worldwide.
\item \textbf{Answer:} C
\\ \textit{Options:} A) \$150,000; B) \$75,000; C) \$35,000; D) \$15,000
\\ \textit{Note:} As of 2020, an annual income of approximately \$35,000 places an individual in the richest one percent globally due to the significant disparity in income distribution and living standards between developed and developing countries.
\item \textbf{Answer:} B
\\ \textit{Options:} A) 43\%; B) 58\%; C) 73\%; D) 88\%
\\ \textit{Note:} The correct answer of 58\% reflects the significant concern among US adults regarding the influence of inherent human biases in the development and implementation of computer programs, highlighting the intersection of technology and social perceptions of fairness.
\end{enumerate}

\section*{True / False}
\begin{enumerate}[label=\arabic*.]
\setcounter{enumi}{5}
\item \textbf{Answer:} False
\\ \textit{Options:} A) True; B) False
\\ \textit{Note:} The correct answer is: Rice University Stadium
\item \textbf{Answer:} True
\\ \textit{Options:} A) True; B) False
\\ \textit{Note:} According to the passage: minor
\item \textbf{Answer:} False
\\ \textit{Options:} A) True; B) False
\\ \textit{Note:} The correct answer is: European continent
\item \textbf{Answer:} False
\\ \textit{Options:} A) True; B) False
\\ \textit{Note:} The correct answer is: evangelisch
\item \textbf{Answer:} True
\\ \textit{Options:} A) True; B) False
\\ \textit{Note:} According to the passage: new papists
\end{enumerate}

\section*{Long Form Response}
\begin{enumerate}[label=\arabic*.]
\setcounter{enumi}{10}
\item \textbf{Answer:} Turkey
\\ \textit{Note:} Expected answer: Turkey
\item \textbf{Answer:} British Dependent Territories
\\ \textit{Note:} Expected answer: British Dependent Territories
\end{enumerate}

\vfill
\noindent\textit{End of Answer Key}
\end{document}